\documentclass[../proyecto.tex]{subfiles}
\graphicspath{{\subfix{../images/}}}
\begin{document}

\section{Resumen}

Este proyecto de tesis es una investigación artística que propone la conceptualización y práctica de popusintesíntesis, entendida como la síntesis popular de sintetizadores, y sus aplicaciones al particular de sintetizadores sonoros.

El neologismo popusintesíntesis es una respuesta a lo propuesto por el artista y empresario estadounidense Peter Blasser, quien define como sintesíntesis su práctica de diseño y fabricación de sintetizadores. Se añade la componente popular para situar esta práctica dentro de la cultura popular, los objetos cotidianos, y la aspiración a ser una práctica masiva de goce sonoro.

En esta tesis se proponen distintas estrategias para la síntesis popular de sintetizadores sonoros, con un fuerte énfasis en herramientas electrónicas y computacionales. La elección de herramientas usadas y propuestas también responden a lo popular, siendo de fuente abierta, de bajo costo, legibles, y de licencias permisivas.

Como marco conceptual definir e incentivar la popusintesíntesis, se incluye la investigación realizada en el campo de la luthería de instrumentos sonoros del siglo XX y XXI, con un énfasis en las infraestructuras técnicas y comerciales que los hicieron posibles, además de las metáforas, historia y contexto cultural que las rodean.

A nivel de artefactos técnicos que acompañan esta tesis, se incluyen diseño y fabricación de máquinas que operan como bloques fundamentales de síntesis sonora, para poder remezclar y con ellos construir distintos sintetizadores, y también proponer derivas y adaptaciones a otros medios como los visuales o textuales, expandiéndose así la popusintesíntesis.

% máximo 1 página

\end{document}